\apendice{Anexo de sostenibilización curricular}  

\section{Introducción}
El concepto de desarrollo sostenible según la definición de la Comisión Brundtland 1987 \cite{CRUE2012}, consiste en la satisfacción de las necesidades actuales sin comprometer la capacidad de las futuras generaciones para atender las suyas.
En un contexto académico, la Conferencia de Rectores de las Universidades Españolas \cite{CRUE2012} propone la incorporación, al ámbito universitario,  de contenidos transversales de sostenibilidad para la formación de profesionales capacitados para comprender y afrontar retos que presenta el efecto de su profesión con la sociedad. Entre los criterios generales para la sostenibilización curricular destacan:
\begin{itemize}
    \item La comprensión crítica de los desafíos ambientales y sociales que van surgiendo a medida que pasa el tiempo.
    \item La capacidad de aplicar soluciones innovadoras teniendo en cuenta el contexto sociocultural en donde nos ubicamos.
    \item Atender a las demandas y propuestas de la sociedad para una mejor visibilización a las minorías.
\end{itemize}

El desarrollo del presente Trabajo de Fin de Grado se han englobado varios aspectos sostenibilidad en el desarrollo de UV-Dermis. La función encargada de solicitar, limpiar y descargar datos meteorológicos optimiza el uso de recursos abiertos de manera eficiente, evitando el colapso del servidor por saturación por repetidas llamadas a la API de la AEMET.
Por otro lado, la plataforma web incluye un diseño accesible de uso intuitivo permitiendo que cualquier persona pueda obtener información preventiva sin necesidad de conocimientos tecnológicos profundos. Asimismo, al ofrecer recomendaciones personalizadas para prevenir y proteger del desarrollo del melanoma maligno de la piel, UV-Dermis favorece al ámbito sanitario dado a que incentiva a cuidarse la piel y de esta manera evitar la alta demanda sanitaria, conllevando a su colapso a largo plazo. 

El conjunto de estas prácticas sostenibles no solo optimizan el uso de recursos y facilitan el acceso preventivo a la población, sino que también ayudarían a reducir la presión sobre el sistema sanitario y aportan una proteccion al entorno ambiental intentando reducir llamadas innecesarias hacia servidores externos, disminuyendo el consumo energético de servidores remotos.